\documentclass[10pt,a4paper]{amsart}
\usepackage[margin=19mm]{geometry}
\usepackage{amsmath,amssymb,mathtools}
\usepackage[scr=rsfs,cal=euler]{mathalfa}
\usepackage{xfrac}
\usepackage[unicode,hidelinks,bookmarks=false]{hyperref}
\newcommand{\ie}{i.e.\ }
\newcommand{\cf}{cf.\ }
\newcommand{\resp}{respectively\ }
\newcommand{\Subsection}{\subsection}
\providecommand{\nobreakdash}{\nobreak\hbox{-}}
\setlength{\emergencystretch}{2em}
\multlinegap0pt
\usepackage{amssymb}
\usepackage{xcolor}
\usepackage{tikz}
\usepackage{enumerate}
\newtheorem{corollary}{Corollary}
\newtheorem{lemma}{Lemma}
\newtheorem{proposition}{Proposition}
\newcommand{\N}{\mathbb{N}}
\newcommand{\R}{\mathbb{R}}
\newcommand{\Rd}{\mathbb{R}^4_\uparrow}
\newcommand{\X}{\times}
\newcommand{\al}{\alpha}
\newcommand{\bx}{\mathbf{x}}
\newcommand{\by}{\mathbf{y}}
\newcommand{\bz}{\mathbf{z}}
\newcommand{\cB}{\mathcal{B}}
\newcommand{\cL}{\mathcal{L}}
\newcommand{\cl}[1]{{\left\lceil #1 \right\rceil}}
\newcommand{\fB}{\mathfrak{B}}
\newcommand{\fl}[1]{{\left\lfloor #1 \right\rfloor}}
\newcommand{\ka}{\kappa}
\newcommand{\la}{\lambda}
\title{The directed landscape from Brownian motion}
\author{Duncan Dauvergne, B\'alint Vir\'ag}
\date{Source: arXiv:2405.00194v4 (CC BY 4.0)}
\begin{document}
\maketitle

\noindent\emph{Source and licence.} The directed landscape from Brownian motion. Version arXiv:2405.00194v4 (CC BY 4.0), distributed by arXiv under the Creative Commons Attribution 4.0 International licence, http://creativecommons.org/licenses/by/4.0/, as recorded in the arXiv licence metadata for this exact version and shown on the abstract page https://arxiv.org/abs/2405.00194v4. Full licence notice: \texttt{licenses/arxiv-2405.00194-CC-BY-4.0.txt}.\par\smallskip

\noindent\textit{Editorial scope.} Counted unit: the article's proposition P:busemann-quadrangle (source0.tex lines 2840--2858). The article's complete original proof of it is delivered with it (source0.tex lines 2904--2966, one \texttt{proof} environment of 63 lines, which the article places away from the statement). No mathematics is written, repaired or re-derived: the statement and the proof are the article's own lines, in the article's own notation, and no retained line is substituted or re-flowed.  Retained uncounted as the source-local context the proof needs: lines 2447--2458; lines 2809--2821; lines 2862--2873; lines 2883--2886.  Nothing was omitted from any retained span, so the package carries no omission record.  No retained line contains a citation, so the wrapper carries no bibliography and no bibliography entry.  No retained line contains a figure environment, an image-inclusion command or drawing code, so no image is delivered and the package declares no asset dependency.  Editorial formatting only, in addition: an AMS \texttt{amsart} preamble that restates the article's own declarations and macros used by the retained lines, namely the environments \texttt{corollary}, \texttt{lemma}, \texttt{proposition} on separate counters, together with the standard packages those lines need.  The retained environments print in this document as Corollary 1, Lemma 1, Proposition 1 in order of appearance, whereas the article prints the same items with the numbering of its own full document.  The complete native source of the article as distributed in the arXiv e-print for this exact version is bundled unchanged as \texttt{sources/158417/source0.tex}.
\begin{corollary}
	\label{C:hat-B-Busemann-landscape}
For $\theta, \bx \in \R^k_\le$ let $\Pi(\theta) = (I_1, \dots, I_{m(\theta)})$ and define
$$
\bar \fB^\theta(\bx, t; \cL) = \lim_{s \to -\infty} \cL(\theta |s|, s ; \bx, 0) - \sum_{j=1}^{m(\theta)}  \cL(s, \theta^{I_j} |s| ; \bx^{I_j}, 0).
$$
Then for any fixed $\theta \in \R^k_\le$, almost surely $\bar \cB^\theta(\bx, t; \cL) = \cB^\theta(\bx, t; \cL)$ for all $\bx \in \R^k_\le, t \in \R$ and
\begin{equation*}
\bar \fB^\theta(\bx, t; \cL) = \lim_{s \to -\infty} \cL(\pi(s), s ; \bx, 0) - \sum_{j=1}^{m(\theta)}  \cL(s, \pi^{I_j}(s) ; \bx^{I_j}, 0)
\end{equation*}
for any disjoint $k$-tuple $\pi$ ending at $(\bx, 1)$ in direction $\theta$. 
\end{corollary}

\begin{lemma}
	\label{L:strong-quadrangle}
	Let $(\bx, s; \by, t) \in \mathfrak X_\uparrow$, and let $\pi = (\pi_1, \dots, \pi_k)$ be the a.s.\ unique optimizer from $(\bx, s)$ to $(\by, t)$. Fix  $m \in \{0, \dots, k\}$, and consider $x_m \le a \le x_{m+1}$ and $y_m \le b \le y_{m+1}$, where we use the convention that $x_0 = y_0 = -\infty$ and $x_{k+1} = y_{k+1} = \infty$. Let 
	\begin{align*}
	R &= \{(z, r) \in \R \X [s, t] : \pi_m(r) < z < \pi_{m+1}(r)\},
	\end{align*}
	and let $\bar R$ be the closure of $R$. Here again we use the convention that $\pi_0 = -\infty, \pi_{k+1} = \infty$. 
	Let $\bx_a = (x_1, \dots, x_m, a, x_{m+1}, \dots, x_k) \in \R^{k+1}$ and similarly define $\by_b$. Then a.s.\
	\begin{align}
	\label{E:xsyt-cha}
	\cL(\bx, s; \by, t) + \cL(a, s; b, t \mid R) \le \cL(\bx_a, s; \by_b, t) \le \cL(\bx, s; \by, t) + \cL(a, s; b, t \mid \bar R).
	\end{align}
\end{lemma}

\begin{lemma}
	\label{L:one-side-bound}
Let $u =(x, s; y, t) \in \Rd$, and let $R \subset [s, t] \times \R$. Let $[s', t'] \subset [s, t]$, and let $\tau:[s', t'] \to \R$ be any geodesic with $\mathfrak g \tau \subset R$. Then we can write
$
\cL(x, s; y, t \mid R) = \sup_\pi \|\pi\|_\cL,
$
where the supremum is over all paths $\pi$ for which the set
$$
\{r \in [s, t] : \pi(r) = \tau(r)\}
$$
is a (possibly empty) closed interval.
\end{lemma}

\begin{lemma}
	\label{L:stays-inside}
Fix $u = (x, s; y, t)$ with $t < 0$ and a point $(z, r)$ with $r < s$ and let $w \in \R$. Then almost surely, the restricted landscape values $\cL(z, r; w, 0 \mid S_u), \cL(z, r; w, 0 \mid S_{y, t}^+)$ are achieved by paths that stay in the interior of the sets $S_u, S_{y, t}^+$ respectively.
\end{lemma}

\begin{proposition}
	\label{P:busemann-quadrangle}
	Consider deterministic sequences $\bx^n, \la^n \in \R^n_\le$ and $\by^n, \kappa^n \in \R^n_\le$. Let $\pi^n$ be the a.s.\ unique semi-infinite optimizer in direction $\la^n$ to $(\bx^n, 0)$ and let $\tau^n$ be the a.s.\ unique semi-infinite optimizer in direction $\ka^n$ to $(\by^n, 0)$. Next, for $x \in \R, t < 0$, let $R_{x, t}^\pm:[0, \infty) \to [-\infty, \infty]$ be the function with $R_{x, t}^\pm(t) = x$ and $R_{x, t}^\pm(r) = \pm \infty$ for $r \ne t$, and suppose that the following conditions hold almost surely:
	\begin{itemize}
		\item For some $u = (x, s; y, t) \in \Rd$ with $t < 0$, we have $\pi^n_n \to_{\mathfrak h} R^-_{x, s}$ and $\tau^n_1 \to_{\mathfrak e} R^+_{y, t}$ as $n \to \infty$.
		\item For any $\al \in \R$, there exists $n_0 \in \N$ such that 
		$
		\pi^n_n(r) < \al r < \tau^n_1(r)  
		$
		for all $r \le s - 1$ and $n \ge n_0$.
	\end{itemize}  
	Then for any $\mu, z \in \R$, the following two convergences hold a.s.
	\begin{align}
	\label{E:slit-convergence}
	\lim_{n \to \infty} \mathfrak B^{(\la^n, \mu, \kappa^n)}(\bx^n, z, \by^n) - \mathfrak B^{(\la^n, \kappa^n)}(\bx^n, \by^n) &= \mathfrak B^{\mu}(z \mid S_u) \\
	\label{E:single-slit-convergence}
	\lim_{n \to \infty} \mathfrak B^{(\mu, \kappa^n)}(z, \by^n) - \mathfrak B^{\kappa^n}(\by^n) &= \mathfrak B^{\mu}(z \mid S_{y, t}^+).
	\end{align}
\end{proposition}

\begin{proof}[Proof of Proposition \ref{P:busemann-quadrangle}]
	We only treat the double-slit convergence \eqref{E:slit-convergence}, as the single-slit convergence \eqref{E:single-slit-convergence} is similar but simpler.
	We first use Lemma \ref{L:strong-quadrangle} to bound the difference $\mathfrak B^{(\la^n, \mu, \kappa^n)}(\bx^n, z, \by^n) - \mathfrak B^{(\la^n, \kappa^n)}(\bx^n, \by^n)$ above and below. 
	
	First, observe that the two bullet points in the lemma guarantee that $x^n_n, \la^n_n \to -\infty$ as $n \to \infty$ and $y^n_1, \kappa^n_1 \to \infty$ as $n \to \infty$. Moreover, these bullets implies that $\pi^n, \tau^n$ are disjoint for large enough $n$, and so $(\pi^n, \tau^n)$ is a semi-infinite disjoint optimizer in direction $(\la^n, \kappa^n)$ to $((\bx^n, \by^n); 0)$. Now let
	$$
	\tilde \pi^n_\mu(r) = (\pi^n(r), \mu |r|, \kappa^n(r)), \qquad \bz^n := (\bx^n, z, \by^n).
	$$
	By Corollary \ref{C:hat-B-Busemann-landscape}, we have that almost surely:
	\begin{equation}
	\label{E:mathfrak-Bla}
	\begin{split}
&\mathfrak B^{(\la^n, \mu, \kappa^n)}(\bx^n, z, \by^n) - \mathfrak B^{(\la^n, \kappa^n)}(\bx^n, \by^n) \\
= &\lim_{r \to -\infty} \cL(\tilde \pi^n_\mu(r), r; \bz^n, 0) - \cL((\pi^n, \kappa^n)(r), r; (\bx^n, \by^n), 0) - \cL(\mu |r|, r; z, 0).
	\end{split}
	\end{equation}
	Next, define
	$$
	R_n := \{(z, r) \in \R \X (-\infty, 0] : \pi^n_n(r) < z < \tau^n_1(r)\}.
	$$
	By Lemma \ref{L:strong-quadrangle}, \eqref{E:mathfrak-Bla} is bounded above by
	\begin{equation}
	\label{E:limsup-above}
	X_n^+ := \limsup_{r \to -\infty} \cL(\mu |r|, r; z, 0 \mid \bar R_n) - \cL(\mu |r|, r; z, 0),
	\end{equation}
	and bounded below by
		\begin{equation}
	\label{E:liminf-below}
	X_n^- := \liminf_{r \to -\infty} \cL(\mu |r|, r; z, 0 \mid R_n) - \cL(\mu |r|, r; z, 0).
	\end{equation}
	To complete the proof, we just need to show that for all large enough $n$:
	\begin{equation}
	\label{E:X-limit}
X_n^- = X_n^+ = \mathfrak B^{\mu}(z \mid S_u).
	\end{equation}
	Now, we can find points $z_- < z < z_+ \in \R$ such that there are leftmost semi-infinite geodesics $\pi_\pm$ in direction $\mu \pm 1$ to $(z_\pm, 0)$ such that
	\begin{equation}
\pi_-(s) < x < \pi_+(s), \qquad  \pi_-(t) < y < \pi_+(t).
	\end{equation}
	Now, let $M > 0$ be large enough so that $\pi^-(r) \le \mu |r| \le \pi^+(r)$ for all $r \le -M$. By the second bullet point, there exists $M' > M$ such that $\pi_\pm|_{(-\infty, -M']} \subset R_n$ for all large enough $n$. Therefore by Lemma \ref{L:one-side-bound}, for large enough $n$, in the definition of 
	$
	\cL(\mu|r|, r; z, 0 \mid A)
	$
	where $A = R_n, \bar R_n, S_u$,
	it suffices to consider paths that stay in the set
	$$
	G := \{(w, r') \in \R \X (-\infty, -M'] : \pi_-(r') \le w \le \pi_+(r')\} \cup (\R\X [-M', 0]).
	$$
	Next, by Lemma \ref{L:stays-inside}, for each $m \in \N$ almost surely we can find paths $\tau_-, \tau_+$ that achieve the restricted landscape values
	$\cL(\fl{\pi_-(-\cl{M'})}, -\cl{M'}; z_-, 0 \mid S_u)$ and 	$\cL(\cl{\pi_+(-\cl{M'})}, -\cl{M'}; z_+, 0 \mid S_u)$ and stay in the interior of the set $S_u$. The first bullet point then implies that $\mathfrak{g} \tau_\pm \subset R_n$ for all large enough $n$. Therefore for large enough $n$ and $r \ge M'$, it suffices to further restrict our paths to the region
	$$
	G' := G \cap (\{(w, r') \in \R \X (-M', 0] : \tau_-(r') \le w \le \tau_+(r')\} \cup \R\X (-\infty, -M'])
	$$
	in the definitions of $
	\cL(\mu|r|, r; z, 0 \mid A)
	$
	where $A = R_n, \bar R_n, S_u$.
	Noting that $G' \subset S_u \cap R_n \cap \bar R_n$ for all large enough $n$, we get that 
	$$
	\cL(\mu|r|, -r; z, 0 \mid R_n) = \cL(\mu|r|, -r; z, 0 \mid \bar R_n)= \cL(\mu|r|, -r; z, 0 \mid S_u)
	$$
	for $r \ge M'$ and all large enough $n$. This yields \eqref{E:X-limit}, completing the proof.
\end{proof}

\end{document}
